(1) 根据法拉第定律与电容压降的约定，回路满足 $L\ddot q_1+M\ddot q_2+q_1/C=0$、$M\ddot q_1+L\ddot q_2+q_2/C=0$。两回路中互感加速度项符号一致，因为磁能交叉项为 $M I_1I_2$，不是额外的耗散阻力。
(2) 引入 $q_+=q_1+q_2$ 与 $q_-=q_1-q_2$，分别有 $(L+M)\ddot q_++q_+/C=0$ 与 $(L-M)\ddot q_-+q_-/C=0$。频率 $\omega_+=1/\sqrt{C(L+M)}$、$\omega_-=1/\sqrt{C(L-M)}$；$M>0$ 时对称模频率较低，互感翻号则两种模序互换。
(3) 初值给 $q_+(0)=q_-(0)=Q$，所以 $q_1(t)=\frac Q2(\cos\omega_+t+\cos\omega_-t)$、$q_2(t)=\frac Q2(\cos\omega_+t-\cos\omega_-t)$。两频率接近时利用和差化积可见，单回路能量在慢尺度上交换；整个无损系统却并无能量衰减。
(4) 总能量 $\mathcal E=[L(\dot q_1^2+\dot q_2^2)+2M\dot q_1\dot q_2]/2+(q_1^2+q_2^2)/(2C)$。微分后代入回路方程得 $\dot{\mathcal E}=0$。电感二次型的特征值 $L\pm M$ 都正，故能量正定；当 $M\to0$ 时 $\omega_+=\omega_-=1/\sqrt{LC}$，因此 $q_2(t)=0$，符合两独立回路的初值。
